% TikZ 图：光滑斜面到粗糙水平面
% Day 29：模型复盘
\begin{tikzpicture}[scale=0.88,>=Latex,line cap=round,line join=round]
  % 左高右低的光滑斜面与粗糙水平面顺接，运动方向与题意一致。
  \coordinate (A) at (0,1.75);
  \coordinate (B) at (3.00,0);
  \coordinate (C) at (7.35,0);
  \draw[thick] (A) -- (B) -- (C);
  \begin{scope}[shift={(1.30,0.99)},rotate=-30.26]
    \draw[fill=gray!12,thick] (-0.32,0) rectangle (0.32,0.42);
  \end{scope}
  \draw[<->,blue,thick] (3.00,0) -- (3.00,1.75)
    node[midway,right] {$h$};
  \draw[->,red,thick] (3.28,0.25) -- (4.26,0.25);
  \node[red,font=\small,fill=white,inner sep=1pt] at (3.77,0.58) {$v$};
  \draw[<->,orange,thick] (3.35,-0.35) -- (6.75,-0.35)
    node[midway,below] {$s$};
  \node[font=\small] at (1.10,2.05) {光滑段：机械能守恒};
  \node[font=\small] at (5.55,0.92) {粗糙段：动能定理};
  \node[font=\small] at (5.30,-0.02) {$\mu$};
\end{tikzpicture}
